\scanpage{90}\mspage{T-4}

3. \uline{Forme canonique} ( $p\ne 2,3$)

C'est la forme
\[
 y^{2}=x^{3}+a_{4}x+a_{6}\qquad\qquad \pi=dx/2y
\]

On a :
\[
 c_{4}=-48\,a_{4}\ ,\quad c_{6}=-4\mathord{.}216\,a_{6}\ ,\quad
 \handwritten{\Delta}=16(-4a_{4}^{3}-27\,a_{6}^{2}).
\]

En posant $Y=2y$, $X=x$ , on retombe sur la forme de Weierstrass
\[
 Y^{2}=4X^{3}-g_{2}X-g_{3}\qquad\qquad \pi=dX/Y
\]
avec
\[
\begin{aligned}
g_{2}&=-4a_{4}=\frac{1}{12}\,c_{4} & g_{3}&=-4a_{6}=\frac{1}{216}\,c_{6}\\
\handwritten{\Delta}&=g_{2}^{3}-27g_{3}^{2} & \unsure{j|J}&=1728\,g_{2}^{3}/\handwritten{\Delta}\ .
\end{aligned}
\]

Les isomorphismes d'une telle courbe sur une autre sont donnés simplement par :
\[
 x\unsure{\circ|\mathord{.}}f=u^{2}x'\qquad\qquad y\;f=u^{3}y'\ .
\]

Deux courbes ayant le même $j$ sont isomorphes sur la clôture séparable de $k$ (corps de base - si c'est un schéma de base, il faut en prendre un revêtement). Plus précisément :

\uline{Cas général} : $j\ne 0,1728$\quad (i.e. $a_{4}\ne 0$, $a_{6}\ne 0$) .

Les deux courbes $C$ et $C'$ sont isomorphes si et seulement si $j=j'$ et $a_{4}a'_{6}/a'_{4}a_{6}$ est un carré ; le plus petit corps sur lequel $C$ et $C'$ deviennent isomorphes est le corps $k'$ obtenu en adjoignant à $k$ la racine carrée de $a_{4}a'_{6}/a'_{4}a_{6}$ .

Les automorphismes de $C$ correspondent à $u=\pm 1$.

\uline{Cas}\quad $j=1728$\quad (i.e. $a_{6}=0$) .

Ici $C\approx C'$ équivaut à $j=j'$ et $a_{4}/a'_{4}\in k^{\ast 4}$ .

Les automorphismes sont donnés par les $u$ tels que $u^{4}=1$.

Courbe type : $y^{2}=x^{3}-x$ .
% archive sheet number: 89

